\chapter{Testing and Validating Compilers}\label{ch:testing-validating-compilers}

A compiler that cannot be checked cannot be trusted.
Optimizations rewrite programs in ways their authors never wrote, and each rewrite can introduce a defect that no single example exposes.
Testing therefore needs structure: contracts that state what each pass preserves, oracles that judge results, and regression suites that remember every past failure.
This chapter builds that structure from pass contracts to minimized failing tests.
We reuse the fault-isolation ideas of the opening chapter, restated here in the language of full pipelines.

\section{Pass Contracts State the Promise}\label{sec:ch37-contracts}

A \term{pass contract} names the input shape a pass accepts, the property it preserves, and the output shape it guarantees.
The input shape lists required invariants: well-formed control flow, valid SSA, or consistent types.
The preserved property names the observable behavior: result values, effect order, or termination on defined inputs.
The output shape lists the invariants the next pass may assume.
A contract makes blame assignable.
When a pipeline miscompiles, the first pass whose output violates its contract is the suspect.
Consider a folding pass over pure expressions.
Its contract accepts well-formed trees, preserves evaluation results on all inputs, and guarantees no new effects.
A test then checks three things: the input satisfied the entry shape, the result matches the oracle, and the output satisfies the exit shape.
Figure~\ref{fig:ch37-pipeline} shows contracts as gates between passes.
Each gate checks its entry shape, runs the pass, and checks the exit shape before handing off.
Table~\ref{tab:ch37-oracles} compares oracle kinds that judge pass results.
Differential testing compares two pipelines or optimization levels on the same input.
Reference interpreters compare compiled behavior against a slow trusted evaluator.
Metamorphic checks compare related inputs whose outputs must satisfy a known relation.
Each oracle catches defects the others miss, so suites combine them.

\begin{figure}[htbp]
\centering
\begin{tikzpicture}[node distance=1.0cm and 1.3cm]
\node[basic block] (p1) {pass A};
\node[basic block, right=of p1] (g1) {contract\\ check};
\node[basic block, right=of g1] (p2) {pass B};
\node[basic block, right=of p2] (g2) {contract\\ check};
\node[basic block, right=of g2] (p3) {pass C};
\draw[cfg edge] (p1) -- (g1);
\draw[cfg edge] (g1) -- (p2);
\draw[cfg edge] (p2) -- (g2);
\draw[cfg edge] (g2) -- (p3);
\end{tikzpicture}
\caption{Contract checks between passes turn a silent miscompilation into a located failure by testing each handoff before the next pass runs.}
\label{fig:ch37-pipeline}
\end{figure}

\begin{table}[htbp]
\centering
\caption{Oracles judge results differently, so a suite combines several kinds rather than relying on one.}
\label{tab:ch37-oracles}
\begin{tabular}{lll}
\toprule
Oracle & Compares & Catches \\
\midrule
Differential & Levels or pipelines & Optimization-only defects \\
Reference & Trusted evaluator & Wrong-code defects \\
Metamorphic & Related inputs & Consistency defects \\
\bottomrule
\end{tabular}
\end{table}

\section{From Failure to Minimal Test}\label{sec:ch37-minimize}

A failing program is usually too large to diagnose.
\term{Test-case reduction} shrinks it while preserving the failure.
The classic loop removes one fragment at a time and keeps the removal when the failure persists.
Algorithm~\ref{alg:ch37-reduce} states the single-fragment greedy version.
It assumes a deterministic failure predicate: the same input always fails or always passes.
It requires that fragments can be removed independently without breaking the test harness.
It returns a test where no single remaining fragment can be removed without losing the failure.
The result is minimal relative to single removals, not globally smallest.
That weaker guarantee is still valuable because each step is cheap and checkable.
The opening chapter applied the same idea to isolate one wrong transformation in a tiny expression language.
The method there built a deterministic reducer, replayed an exact reduction trace, and argued minimality step by step.
We generalize that discipline here to full pipelines.
First, freeze the input, the flags, and the environment so the failure reproduces on demand.
Second, reduce the program while the oracle still reports the same defect signature.
Third, reduce the flags while the same signature persists.
Fourth, record the minimal case with its contract check, trace, and expected output.
A frequent mistake reduces the program and the flags together.
Coupled reduction confuses which change mattered and produces cases that fail for a different reason than the original.
Another mistake accepts flaky failures during reduction.
Nondeterminism makes every removal decision suspect.
Repeated reproduction is a practical filter for flaky failures, not a proof of determinism.
The minimality claim below assumes a deterministic predicate that recognizes the same failure signature, including any required validity checks.

\begin{algorithm}[htbp]
\caption{Greedy single-fragment test-case reduction}\label{alg:ch37-reduce}
\begin{pseudocode}
\Require failing test $T$, finite removable fragments, deterministic same-failure predicate $F$
\Ensure $T$ still fails and no permitted valid single-fragment removal preserves that failure
\State confirm $F(T)$; abort reduction if the original signature is absent
\Repeat
\State $changed \gets false$
\ForAll{$i$ from $1$ to $n$}
\State $T' \gets T$ with $f_i$ removed
\If{$T'$ is well-formed and $F(T')$ holds}
\State $T \gets T'$
\State $changed \gets true$; refresh fragment numbering and restart the scan
\State \Break
\EndIf
\EndFor
\Until{not $changed$}
\State \Return $T$
\end{pseudocode}
\end{algorithm}

\section{Worked Example: Shrinking a Folding Failure}\label{sec:ch37-worked}

We reduce a failing constant-folding case by hand.
The original test folds \code{(x + 7) + (3 + 5)} with $x = 10$ and expects $25$.
The pipeline reports $17$, a wrong-code signature.
Fragments are the outer add, the \code{x + 7} term, the \code{3 + 5} term, and the binding $x = 10$.
We confirm the failure three times before touching anything.
Try removing the outer add: the program no longer parses, so the removal is rejected.
Try removing \code{x + 7}: the remainder \code{3 + 5} folds to $8$, which matches the oracle, so the failure vanishes and the removal is rejected.
Try removing \code{3 + 5}: the remainder \code{x + 7} with $x = 10$ gives $17$ both before and after folding, so no failure appears and the removal is rejected.
Try removing the binding: $x$ becomes undefined and the harness reports a different signature, so the removal is rejected.
No single removal preserves the wrong-code signature.
The test is therefore irreducible under these permitted single-fragment removals, assuming the stated failure predicate; this is not a global minimum-size claim.
Hand arithmetic confirms the diagnosis: $x + 7 = 17$, $3 + 5 = 8$, and $17 + 8 = 25$.
Reporting $17$ means the folder kept the left term and dropped the right term.
The listing below pins the minimal case for the regression suite.
\begin{irlisting}
; minimal folding failure, expect 25
; x = 10
; expr = (x + 7) + (3 + 5)
; observed after fold: 17
; diagnosis: right term dropped
\end{irlisting}
The suite stores the input, the flags, the oracle output, and the contract check that locates the faulty handoff.

\section{Regression Suites That Remember}\label{sec:ch37-regression}

A \term{regression test} locks a past failure so it can never return silently.
Each entry records the minimal program, the exact flags, the environment, the oracle verdict, and the contract check that locates the defect.
Entries run on every change to the passes they cover.
A new failure must either match a known entry or earn a new one through reduction.
\begin{ccnote}[One defect per entry]
An entry that covers two defects passes when either is fixed and hides the other.
Split combined failures before storing them.
\end{ccnote}
We organize entries by the handoff they guard.
Folding entries pin expression contracts.
Frame entries pin calling conventions.
Guard entries pin speculation exits with injected failures.
The \enquote{guard the handoff} rule keeps suites diagnostic: a red entry names the suspect pass directly.
Randomized testing complements the suite.
Generators produce well-formed programs, differential oracles compare levels, and every new mismatch enters reduction immediately.
Coverage guides generation toward untested opcodes and shapes without replacing directed tests.

\section{Generators and Fuzzing Discipline}\label{sec:ch37-fuzzing}

Hand-written tests cover known defects, but they rarely surprise the compiler.
\term{Randomized testing} generates well-formed programs, runs oracles over them, and reduces every mismatch.
Generation needs the same contract discipline as the rest of the suite.
The generator produces only inputs that satisfy the entry shape of the handoff under test.
Anything else tests the validator instead of the pass.
A practical generator grows programs from a small grammar of expressions, control flow, and calls.
It tracks defined names, live values, and effect order so outputs stay meaningful.
It biases toward recently changed opcodes and historically buggy shapes without abandoning the base distribution.
Seeds make every run reproducible.
A failing seed replays exactly, reduces through Algorithm~\ref{alg:ch37-reduce}, and enters the regression suite as a directed entry.
We quarantine flaky seeds rather than deleting them.
A seed that fails once in fifty runs still names a real defect or a real nondeterminism source.
Consider a generator that emits ten thousand straight-line programs per night.
Differential comparison of two optimization levels flags eleven mismatches.
Reduction shrinks them to three distinct minimal cases, and two match known entries.
The night produced exactly one new regression test and one updated coverage gap.
That arithmetic is the point: fuzzing is measured by distinct minimal defects, not by raw program count.
The \enquote{count distinct defects} rule keeps fuzzing honest about its yield.

\section{Implementation Notes and Complexity}\label{sec:ch37-implementation}

Contract checks cost $O(N)$ per handoff where $N$ is the program size, using validators that walk each structure once.
Reduction with $n$ fragments costs $O(n)$ oracle runs per pass over the fragments, each run costing a full pipeline execution.
Restarting after each successful removal permits $O(n^2)$ oracle calls overall when each accepted step strictly reduces the fragment count.
This is single-deletion irreducibility relative to the permitted edits, not globally minimum program size.
Small fragments keep $n$ large but each run cheap; large fragments reverse the trade.
Caching well-formedness results avoids rechecking unchanged regions.
Production suites shard by handoff, run deterministically under fixed seeds, and quarantine flaky entries instead of deleting them.
They version expected outputs with the contract that produced them.
A common pitfall stores outputs without the flags that produced them.
A later flag change then turns the whole suite red for no local reason.
Another pitfall lets generators drift from the language definition.
Generated programs that violate entry shapes test the validator, not the pass.
We freeze grammars with the contracts they target and review generator changes as contract changes.
Every suite entry is a contract with a replayable witness.
% TODO(cite): need a primary reference for delta-debugging minimality guarantees.

\section{Exercises}\label{sec:ch37-exercises}

\begin{enumerate}
\item Run Algorithm~\ref{alg:ch37-reduce} on a five-fragment test where only fragments $2$ and $4$ together trigger failure, list which removals are kept, and state whether the result is globally minimal.\label{ex:ch37-reduce-trace}
\item Write a pass contract for a dead-code pass over straight-line code, naming its input shape, preserved property, output shape, and two tests that each violate a different clause.\label{ex:ch37-contract-write}
\item Design a regression entry for a miscompiled call sequence, listing the program, flags, environment, oracle verdict, and the contract check that locates the handoff.\label{ex:ch37-entry-design}
\end{enumerate}

\section{Further Reading}\label{sec:ch37-reading}

Appel discusses testing compiler components against their specifications~\cite{appel-1998-modern}.
Morgan treats pass structure and the obligations between passes~\cite{morgan-1998-optimizing}.
Safonov examines validation as part of trustworthy compilation~\cite{safonov-2010-trustworthy}.
The C standard draft supplies the conformance rules that C-facing tests must encode~\cite{wg14-2011-c11}.
